% !TEX root = main.tex
% Section 02: Milestones.
% Source of truth: docs/MILESTONES.md.

\chapter{Milestones}

\section{Stage taxonomy}

The lab is organized into six stages. Each stage ends with a
measurable result. The next stage begins only after the previous
stage's result is recorded.

\begin{longtable}{p{1.0cm}p{3.2cm}p{9.0cm}}
\toprule
\textbf{Stage} & \textbf{Title} & \textbf{Goal} \\
\midrule
0 & Lab scaffold & Fork, document, snapshot, governance, CI. \\
1 & Local inference & Tiny SLM behind the existing Node bridge. \\
2 & Cognitive router & Tiny model is the first-stage router. \\
3 & Core/client separation & Bridge becomes \texttt{src/core}. \\
4 & Desktop shell & Tauri-style client around the new Core. \\
5 & Hardware profiler & Adaptive model selection. \\
\bottomrule
\end{longtable}

\section{Stage 0 --- Lab scaffold}

Goal: establish a clean derivative repository of
ResonantOS/2.0.0-alpha with the vision, architecture, governance,
and CI baseline.

Deliverables (in this order):
\begin{itemize}
  \item \texttt{README.md}, \texttt{docs/VISION.md},
        \texttt{docs/ARCHITECTURE.md}, \texttt{docs/MILESTONES.md},
        \texttt{docs/WORKFLOW.md}.
  \item Snapshot of upstream \texttt{dev} in \texttt{baseline/}.
  \item \texttt{.gitignore}, \texttt{LICENSE}, \texttt{CONTRIBUTING.md},
        \texttt{CODE\_OF\_CONDUCT.md}, \texttt{SECURITY.md}.
  \item Lab-flavored Gitflow with \texttt{main} and \texttt{develop}.
  \item Initial decisions recorded in \texttt{docs/decisions/}.
  \item LaTeX documentation layer (this document).
  \item CI workflows: \texttt{pr-checks}, \texttt{upstream-deterministic},
        \texttt{gitflow}, \texttt{latex-build}.
\end{itemize}

\section{Stage 1 --- Local inference}

Goal: prove that a tiny SLM can run on the existing Node bridge
without destabilizing the upstream ResonantOS runtime.

Deliverables:
\begin{itemize}
  \item A local-model service backed by a defensible runtime
        (initially \texttt{bitnet.cpp} or a Rust-native equivalent).
  \item A capability-scoped bridge route exposing it.
  \item Throughput, RAM, and CPU measurements on a reference machine.
  \item A documentation note describing the integration points with
        the upstream Node bridge.
\end{itemize}

Decision gate: continue to Stage 2 only if the local model can run
reliably below 4\,GB RAM on a CPU-only machine and produces structured
output that the bridge can validate.

\section{Stage 2 --- Cognitive router}

Goal: make the tiny local model the first-stage intent/router
before any cloud call.

Deliverables:
\begin{itemize}
  \item A new bridge route that accepts user intent, runs the local
        model, returns a typed capability request.
  \item A deterministic Core that validates the request and either
        executes locally, escalates, or returns a missing-capability
        error.
  \item A benchmark comparing routing accuracy and cloud-call
        reduction against the current ResonantOS provider-routing
        path.
\end{itemize}

Decision gate: continue to Stage 3 only if the local router achieves
the routing-accuracy threshold (set before the stage begins) and
demonstrably reduces cloud calls for common tasks.

\section{Stage 3 --- Core/client separation}

Goal: make the Core addressable by clients other than Chrome.

Deliverables:
\begin{itemize}
  \item A new \texttt{src/core} module that exposes the same
        capability surface as the current Node bridge.
  \item The Chrome extension becomes a thin client talking to the
        Core through the existing authenticated bridge protocol.
  \item A minimal CLI client that issues a single capability
        request end-to-end.
\end{itemize}

Decision gate: continue to Stage 4 only if the CLI client can
perform at least one meaningful task (e.g.\ read a file, summarize
it, write a result) without the Chrome extension loaded.

\section{Stage 4 --- Desktop shell}

Goal: bring back a native desktop client around the new Core.

Deliverables:
\begin{itemize}
  \item A Tauri-style shell that loads a web UI and talks to the
        Core.
  \item macOS \texttt{.app} and \texttt{.dmg} builds.
  \item Windows \texttt{.msi} and \texttt{.exe} builds.
  \item Linux package builds.
\end{itemize}

Decision gate: continue to Stage 5 only if both macOS and Windows
installers build successfully on the reference CI and pass a
one-click smoke test.

\section{Stage 5 --- Hardware profiler}

Goal: detect the host hardware and choose the model tier adaptively.

Deliverables:
\begin{itemize}
  \item A hardware-profiler service that reports architecture, CPU,
        RAM, GPU/NPU, disk, and OS.
  \item A model-tier selector that maps the profile to a recommended
        Cognitive Core model and provider set.
  \item End-to-end benchmarks across at least three reference
        machines: low-RAM, mid-range, accelerated.
\end{itemize}

\section{Working assumptions}

\begin{itemize}
  \item The local SLM is honest about uncertainty. It returns
        structured fields, not free-form prose, for routing
        decisions.
  \item The deterministic Core never blocks on the local model for
        more than a bounded budget. If the model times out, the Core
        falls back to a safe default.
  \item Capability policy is written before the model is integrated,
        not after.
  \item All measurements are reproducible from a fresh checkout.
\end{itemize}
